% Start preamble
\documentclass[12pt,a4paper]{article}
\usepackage{geometry}
 \geometry{
 a4paper,
 total={170mm,257mm},
 left=20mm,
 top=20mm,
 }
\usepackage{currfile}
\usepackage[utf8]{inputenc}
\usepackage[T1]{fontenc}
\usepackage[pdftex]{graphicx}
\graphicspath{{./}}
\usepackage{enumitem}
\usepackage{pdfpages}
\usepackage{hyperref}
\usepackage{tikz}
\usepackage{attachfile}
\usepackage{epstopdf}
\usepackage{array}
\usepackage{multirow}
\usepackage{multicol}
\usepackage{float}
%\usepackage[table]{xcolor,colorbl}
\setlength{\textwidth}{16cm}
\setlength{\oddsidemargin}{-0.5cm}
\setlength{\evensidemargin}{-0.5cm}
%\setlenght{\headsep}{0cm}
\setlength\parindent{0pt}
%\setlength{\extrarowheight}{3pt}
\usepackage{listings}
%\usepackage{xcolor}
\input{arduinoLanguage.tex}
%%%%%% Counting oppgaves %%%%%%
 \newcount\questnum \questnum=0
 \def\oppgave{
		}
% End preamble
\begin{document}
\title{TIF Instrumentering - Oppgaver}
\author{Faglærer: Fred-Olav Mosdal 90507684\\}
\maketitle
\oppgave{}%1
\vskip 2.5pt 
Forklar de ulike aggregattilstandene, bruk vann som eksempel.
\vskip 5pt 
\begin{tikzpicture}
	\draw[step=0.5cm,gray!20,very thin]  grid (17,4) ;
\end{tikzpicture}
\vskip 2.5pt 
\oppgave{}%2
\vskip 2.5pt 
Hva menes med en faseovergang?
\vskip 5pt 
\begin{tikzpicture}
	\draw[step=0.5cm,gray!20,very thin]  grid (17,10) ;
\end{tikzpicture}
\vskip 2.5pt 
\oppgave{}%3
\vskip 2.5pt 
Hva er temperatur i grader celsius ved det absolutte nullpunkt?
\vskip 5pt 
\begin{tikzpicture}
	\draw[step=0.5cm,gray!20,very thin]  grid (17,10) ;
\end{tikzpicture}
\vskip 2.5pt 
\oppgave{}%4
\vskip 2.5pt 
Hva er temperaturen i kelvin ved det absolutte nullpunkt?
\vskip 5pt 
\begin{tikzpicture}
	\draw[step=0.5cm,gray!20,very thin]  grid (17,6) ;
\end{tikzpicture}
\vskip 2.5pt 
\oppgave{}%5
\vskip 2.5pt 
Forklar hvordan trykk, temperatur og volum i en gass henger sammen.
\oppgave{}%4
\vskip 2.5pt 
Hva er temperaturen i kelvin ved det absolutte nullpunkt?
\vskip 5pt 
\begin{tikzpicture}
	\draw[step=0.5cm,gray!20,very thin]  grid (17,6) ;
\end{tikzpicture}
\vskip 2.5pt 
\vskip 5pt 
\begin{tikzpicture}
	\draw[step=0.5cm,gray!20,very thin]  grid (17,6) ;
\end{tikzpicture}
\vskip 2.5pt 
\oppgave{}%6
\vskip 2.5pt 
Forklar kontinuitetsloven med et praktisk eksempel.
\vskip 5pt 
\begin{tikzpicture}
	\draw[step=0.5cm,gray!20,very thin]  grid (17,6) ;
\end{tikzpicture}
\vskip 2.5pt 
\oppgave{}%7
\vskip 2.5pt 
Kelvin og celsius:

\begin{itemize}
	\item Hva er 100 *C i kelvin?
\begin{tikzpicture}
	\draw[step=0.5cm,gray!20,very thin]  grid (16,6) ;
\end{tikzpicture}
	\item Hva er 100 K i grader celsius?
\begin{tikzpicture}
	\draw[step=0.5cm,gray!20,very thin]  grid (16,6) ;
\end{tikzpicture}

\end{itemize}
\vskip 5pt 
\begin{tikzpicture}
	\draw[step=0.5cm,gray!20,very thin]  grid (17,4) ;
\end{tikzpicture}
\vskip 2.5pt 
\oppgave{}%8
\vskip 2.5pt 
Regn ut trykket i pascal fra et lodd som virker med 980 N på et areal som er 0,1 m². En dag er det 7 °C, og regnet høljer ned. Dagen etter blir det -10 °C, og alt vannet fryser til is. Så er det minusgrader i flere uker, men allikevel forsvinner isen gradvis. Hvilke aggregattilstander er vannet innom, om vi tar utgangspunkt i regnværet? Sett opp riktig rekkefølge, og finn ut hva de ulike faseovergangene kalles.

\begin{itemize}
	\item I hvilken beholder er trykket høyest?
\begin{tikzpicture}
	\draw[step=0.5cm,gray!20,very thin]  grid (16,6) ;
\end{tikzpicture}
	\item I hvilken beholder er trykket lavest?
\begin{tikzpicture}
	\draw[step=0.5cm,gray!20,very thin]  grid (16,6) ;
\end{tikzpicture}

\end{itemize}
\vskip 5pt 
\begin{tikzpicture}
	\draw[step=0.5cm,gray!20,very thin]  grid (17,6) ;
\end{tikzpicture}
\vskip 2.5pt 
\oppgave{}%9
\vskip 2.5pt 
Nedenfor ser du tre lukkede beholdere. Alle beholderne har samme antall partikler/molekyler, og de er av samme type. Temperaturen er ulik, den er lavest i beholder A og høyest i beholder C.

$$\includegraphics[width=0.4\textwidth]{../output/nogpl/beholdere.png}$$
\vskip 5pt 
\begin{tikzpicture}
	\draw[step=0.5cm,gray!20,very thin]  grid (17,6) ;
\end{tikzpicture}
\vskip 2.5pt 
\oppgave{}%10
\vskip 2.5pt 
På figuren ser du en trykkmåler som er plassert i bunnen av en tank. Her er Pa trykket i toppen av tanken, p er tettheten til væsken, og h er væskehøyden målt fra bunnen av tanken. Hva er trykket i bunnen av tanken? Er det høyere, lavere eller det samme som i toppen av tanken?\\

$$\includegraphics[width=0.4\textwidth]{../output/nogpl/trukkitank.png}$$
\vskip 5pt 
\begin{tikzpicture}
	\draw[step=0.5cm,gray!20,very thin]  grid (17,6) ;
\end{tikzpicture}
\vskip 2.5pt 

\end{document}
